\section{The law of gate substitution: statement, five component laws, six path hypotheses}
\label{sec:law}

\noindent\fbox{\begin{minipage}{\dimexpr\linewidth-2\fboxsep-2\fboxrule\relax}
\textbf{The law of gate substitution.} Every gate-centered work package that an
enterprise places in its DGF, whatever gates it has chosen, obtains, and moves to, an effective automated implementation of its accepted input-output
contract. Inquiry, exceptional-case analysis, verification, the general gate's arbitration, and the commitment it
records are included.
The recurring work needed to operate and verify those implementations is itself a substitution
target, not a permanently human role. At fixed governed output, the human workforce required to
execute the DGF contracts toward its remaining support floor, and then toward zero.

\smallskip\emph{Status: a conjecture proposed by the author, not an established result. Its ten-year
cohort form (Section~\ref{sec:testing}) can be refuted.}
\end{minipage}}

\medskip
``Every gate'' is a universal claim over gate types and their declared case domains. It is not a claim
that a machine will know an unobserved fact, that every judgment is error-free, or that all enterprises
cross the boundary on the same date. The accepted quality specification stays fixed while
implementations are compared; the thesis does not make the specification easier by removing difficult
cases. Its technical and adoption components are separately observable: a qualifying implementation
exists, and the organization actually uses it under a valid mandate.

\subsection{Five component laws}

The five laws name distinct commitments of the thesis.

\law{Law 1}{Contract substitutability}{A DGF gate's function is specified by the transformation it
performs, not by the profession that currently staffs it.} Operational contracts, augmented with
inquiry and action interfaces, capture enough of that function to support an implementation-independent
evaluation. It fails for a stated domain if a material part of successful execution cannot be
represented or evaluated without an indispensable personal intervention.

\law{Law 2}{Complete execution}{No DGF gate class is a permanent exception to autonomous execution.}
Progressively broader implementations ultimately cover the entire gate, rather than stopping at a
permanently human layer of difficult decisions. A durable technical or institutional requirement for
human execution in one gate class contradicts its strong form.

\law{Law 3}{Route preservation}{Replacing the execution of every component replaces the execution of
the route without erasing its required gates.} Replacements preserve the composition the enterprise has decided: its gates, their order on each route,
their conditions, and their returns. Gates are interchangeable for the enterprise that composes the
DGF, which can add, remove, merge, or reorder them at any time; they are not interchangeable for the
agent that executes it. What is preserved is every obligation and the logical order of authorizations,
not the physical existence of each box: one service can execute several contracts, but dropping or
reordering a required control to make automation easier is not its replacement.

\law{Law 4}{Execution displacement}{For a fixed amount of governed work, automating the complete route
removes the labor demand that staffed its execution.} The reduction is not merely a change in job
descriptions. It follows when removed execution work is not replaced by equally large required human
support (Section~\ref{sec:workforce}).

\law{Law 5}{Supervisory closure}{Monitoring, exception review, policy maintenance, and verification are
additional transformations to automate, not permanent human refuges assumed by definition.} If they
remain labor-intensive, the system is only partially closed.

\subsection{What makes the law stronger than an augmentation claim}

An augmentation claim permits the same human team to remain indispensable while using better tools.
Complete substitution does not. A staff member who must interpret every output, reconstruct every rare
incident, or sign every individual approval is still executing part of the gate. Reducing the average
time of that intervention is progress toward substitution, not completion of it. Likewise, the future
workload is not assumed to consist only of easy cases. When an exception becomes a named subprocedure,
the thesis treats that subprocedure as a new target for the same input-output analysis. An
indefinitely expanding collection of indispensable human subprocedures would defeat the strong
endpoint, even if each software release automates more standard cases.

\subsection{Six hypotheses about the path}
\label{sec:hypotheses}

The component laws state the destination. Six hypotheses describe the transition and can be tested
long before the endpoint. Each predicts an observable pattern and can fail within a specified case
population. Automation is evaluated against independently assessed quality, not merely whether a
workflow reaches an approval state.

\law{H1}{Gate readiness}{Gate-level readiness improves out-of-sample prediction of reliable automation
beyond model capability and case length.} Accessible inputs, explicit criteria, independent evidence
checks, scoped actions, and testable outcomes explain differences between otherwise comparable gate
implementations (Appendix~\ref{app:readiness}).
\failure{Failure criterion:} prospectively scored readiness provides no reproducible predictive
improvement over case-size, duration, and evidence-completeness baselines, or repeatedly assigns high
readiness to unresolvable critical failures.

\law{H2}{Handoff propagation}{2a: Evidence-bearing, structured handoffs reduce total human work across
the connected process.} Production, verification, reconciliation, and repeat work all enter the
endpoint; the predicted format contrast is $\beta^{(h)}_{F}<0$ for total hours
(Section~\ref{sec:testing}). \textbf{2b: At equivalent downstream benefit and quality, agents produce
and verify the structured handoff at lower resource cost than humans.} The agent contrast is
$\beta^{(c)}_{A}+\beta^{(c)}_{AF}<0$ for upstream production-and-verification cost. This is the
specifically agentic increment, separate from the value of structure.
\failure{Failure criteria:} 2a fails without lower total hours; 2b fails without lower upstream cost
and equivalent downstream benefit and quality. Both comparisons retain all assigned cases, including
rejected artifacts and work transferred upstream.

\law{H3}{Labor concentration in the human residual}{As agents absorb the standard path, residual cases
retain a larger fraction of baseline labor than their fraction of case volume.} A small human queue
can therefore contain a large share of the work. The prediction concerns baseline effort measured
before deployment, not a difficulty score defined after observing the routing decision.
\failure{Failure criterion:} the residual does not become enriched in baseline effort under the tested
standard-case rollout, or this pattern disappears across held-out case classes. Random-audit and
mandatory-signature queues provide comparison groups.

\law{H4}{Residual workload prediction}{Exception-adjusted forecasts predict human-capacity change
better than the uniform-effort model $q+(1-q)m$.} The comparator gives every exception the baseline
average effort and sets review to fraction $m$ of that average. The component model additionally represents baseline selection, post-deployment exception effort,
rework, and recurring human upkeep. Beating a comparator that omits hours actually worked would prove
little, so a second comparator receives the same rework and upkeep hours and the directly estimated
conditional means $h_E$ and $h_R$. With coherent inputs it is algebraically identical to the component
model; the separation of $\kappa$ from $\eta$ earns its keep only when the routing rule, the case mix,
or the agent configuration changes, because $\phi$ can then be re-estimated on historical work while
$\eta$ and $\mu$ are carried over. H4 therefore tests a transferability hypothesis: that effort ratios
carry over to new conditions better than conditional hours do. It is plausible where assistance scales
with the baseline difficulty of a case, and doubtful where it helps some categories of exception much
more than others, since a change of mix then changes the average $\eta$ without any change in the
agent.
\failure{Failure criterion:} using prospectively estimated inputs, the component model does not
improve held-out total-hours forecasts over the uniform comparator and a stratified case-mix baseline,
does not beat the matched comparator after a change of routing, mix, or configuration, or exceeds a
prespecified error tolerance.

\law{H5}{Delegation before authorization}{At consequential gates, delegation of analytical production
advances ahead of delegation of final authorization.} The enterprise changes who prepares and executes
a decision before changing who is empowered to approve it and answer for its consequences. H5
describes the \emph{order} of the transition; Law~2 states its endpoint.
\failure{Failure criterion:} analysis and authorization transfer at the same rate, or authorization
transfers earlier, in settings where both can vary. A universal human-signature mandate explains a
fixed boundary but cannot by itself validate the prediction. The lag is an institutional outcome to
measure, not a claim about a uniquely human cognitive capability \citep{novelli2024}.

\law{H6}{Residual exhaustion}{The set of indispensable human interventions shrinks toward empty:
categories of intervention close faster than new ones appear, including those created by supervising
the agents.} H1--H5 can all hold while a small team remains indispensable. Three propositions of increasing strength
must not be confused: one human intervention can be automated; the remaining interventions become
automatable one after another; and there is a date after which no necessary human intervention remains.
The law needs the third. H6 is the hypothesis that separates complete substitution from durable
complementarity. Each remaining intervention is recorded
by category and by origin: present at baseline, displaced from another role or organization, or created
by the automation itself.
\failure{Failure criterion:} over three successive annual evaluations, the number of indispensable
intervention categories, or the hours they absorb, stops falling, or newly created categories offset those that close. Such a result rejects this trajectory
hypothesis; like the plateau of Section~\ref{sec:testing}, it is warning evidence for the endpoint, not
its refutation before the deadline.

\section{Three worked routes and their substitutable transformations}
\label{sec:routes-worked}

Each generic route is illustrated by one case from the practitioner framework. The initial situations
and the recorded opinions come from that framework; the route orders are the illustrative configuration
of Figure~\ref{fig:workflows} (Section~\ref{sec:notfixed}), and the proposed machine operations are the
implementation examined by the thesis. The cases are not reports of tested autonomous deployments.

\subsection{W1, Buy: a new platform, from supplier request to the general gate}

The framework's SaaS CRM serves 2,000 sales users and handles confidential customer data. It plans SSO
with optional MFA, lacks contractual SIEM log export, and has untested restoration with no defined
recovery-time objective. Its Security opinion is red, with log export to be negotiated before
signature and MFA and recovery conditions before go-live \citep[p.~103]{canale2026}. The same dossier
gives work to every gate: Procurement and Legal carry the negotiation of the log-export clause,
Compliance qualifies the customer data, IT checks the catalog fit and the run model, Finance validates the
cost, and the general gate arbitrates. The
Security opinion is one transformation among them.

\begin{table}[!htb]
\centering\small
\caption{W1, Buy, in the illustrative order. The proposed machine operations preserve the CRM case's
initial red opinion; they do not assert a later approval.}
\label{tab:w1}
\begin{tabularx}{\linewidth}{@{}P{2.7cm}LL@{}}
\toprule
Gate & Input & Output and automation target\\
\midrule
Procurement & Need, budget, requested controls, supplier proposals and evidence. & Sourcing questions,
offer comparison, evidence requests, due-diligence result.\\
Legal & Draft contract, service-level and data annexes. & Redlined clauses: log export, audit right,
liability, reversibility.\\
Compliance & Data categories, purposes, hosting locations. & Processing assessment, transfer and
retention conditions.\\
Security & Supplier security package, flows, identity model, classification. & Evidence-backed opinion;
contractual changes and preconditions, including the MFA and recovery conditions.\\
IT & Catalog, run and support model. & Catalog fit, run and support model, operating constraints.\\
General & Opinions, reservations, business case, budget, owners. & Go, go with reservations, rework, or
no-go, with owners and dates.\\
\bottomrule
\end{tabularx}
\end{table}

In this configuration the enterprise consults Security before the run model is fixed; automation must respect the order the enterprise has chosen and must not send a
technically integrated platform to Security after the fact. The reusable object is the evidence record of
Section~\ref{sec:stages}, here a contract with exact clause locators; Legal, Compliance, and Security
reuse it instead of re-extracting the same PDF. An automated red disposition is a \emph{successful}
execution of this example when the reference requirements still fail. Automatically producing a green
status would not be substitution; it would be degradation of the governance function. The supplier
belongs to another organization: buyer-side automation cannot hide the supplier's evidence-production
labor, and the complete claim for this route includes that interaction.

\subsection{W2, Integrate: an information system inherited from an acquisition}
\label{sec:cases-parameters}

The acquired entity has 400 people, its own on-premises Active Directory, cloud mail, and a
provider-hosted ERP. A payroll requirement motivates network interconnection within thirty days,
despite missing account inventories, no endpoint detection and response (EDR), local logs, and an
undocumented past incident \citep[p.~104]{canale2026}. IT first frames integration, isolation, or
retirement; Architecture and Security then assess that scope, Legal and Procurement recover the
inherited ERP obligations, Compliance qualifies the inherited data, the general gate's support staff
track the dated obligations, and Finance prices the integration options. The corporate acquisition
decision itself is outside this route.

\begin{table}[!htb]
\centering\small
\caption{W2, Integrate, in the illustrative order.}
\label{tab:w2}
\begin{tabularx}{\linewidth}{@{}P{2.7cm}LL@{}}
\toprule
Gate & Evidence received & Work, output, and automation target\\
\midrule
IT & Payroll need; inherited AD, mail, ERP; available inventories and provider documents. & Frame what
will be integrated, isolated, or retired; plan the run takeover; transmit scope, dependencies, and
explicit evidence gaps.\\
Architecture & IT scope, inherited flows and interfaces, network constraints. & Target integration
architecture; design of the restricted connection; conditions.\\
Security architecture & Scope; identity, exposure, incident, EDR, and logging evidence. & Assess
integration risk; carry the framework's conditional (yellow) opinion into a risk card, a restricted
connection scope, and dated control requirements.\\
Legal & Hosted ERP contract, inherited service obligations. & Obligations register; novation or exit
conditions.\\
Compliance & Inherited personal data, applicable rules. & Conditions on data access and retention.\\
General & Gate opinions, residual risks, integration owner, cost and schedule. & Record the decision and
its authorized scope, accepted residual risk, conditions, and follow-up; preserve the integration
director's risk ownership.\\
\bottomrule
\end{tabularx}
\end{table}

The framework's path overview prohibits interconnection before the acquisition audit. Its worked case
records a bounded exception: a filtered partner-zone connection during the audit without joining
directories, EDR before any connection, the account inventory within thirty days, SIEM log integration
within sixty days, and a risk card signed by the integration director; directories are joined only
after the audit \citep[pp.~18, 104]{canale2026}. These are different prerequisites and dated
obligations, not a blanket approval.

An agent can reconcile inventory exports, link accounts to owners, locate provider reports, and
maintain the reservations register that later gates consume instead of reconstructing the audit. The
model of Section~\ref{sec:frontier} makes the remaining work explicit: investigating the undocumented
incident concentrates baseline effort in exceptions ($\kappa$, hence $\phi$); pursuing unavailable
provider reports determines how much effort persists after assistance ($\eta$); recurring condition
checks consume standard-path review ($\mu$); and maintaining integration playbooks consumes fixed
upkeep ($B_A$). A missing report remains missing: an agent must obtain evidence, run an accepted
alternative investigation, or issue a justified unresolved or refusal disposition. A complete-looking
record does not complete an audit. Low repetition, inherited state, and supplier dependence make
Integrate the most demanding test of the law; shared provision addresses scale, not those information
deficits.

\subsection{W3, Build: a new internal project}

The third case places the framework's internal IT-support agent on the Build route
\citep[pp.~19, 105]{canale2026}. The proposed service uses Copilot Studio and Microsoft Foundry
(formerly Azure AI Foundry; \citealp{msfoundry2026}), reads tickets and SharePoint through its
integrations, and accesses Microsoft Graph. It uses the developer's identity and service principal, can
reset passwords without approval, lacks prompt and action logs, and retrieves across the whole
SharePoint index. Tickets contain personal data. The framework rejects this design with a red opinion
and proposes a read-only knowledge-base pilot without Graph as an alternative. Around that opinion, IT
checks the catalog and the run model, architects frame the target design, readiness experts review
tests and runbooks, the general gate's support staff and Finance assemble schedule and budget, and
sponsors commit at the general gate.

\begin{table}[!htb]
\centering\small
\caption{W3, Build, in the illustrative order. Rows after the Security opinion describe required reviews
of a revised design, not decisions already obtained.}
\label{tab:w3}
\begin{tabularx}{\linewidth}{@{}P{2.7cm}LL@{}}
\toprule
Gate & Evidence received & Work, output, and automation target\\
\midrule
IT & Support need; proposed service and platforms, costs, operating constraints. & Catalog conformity,
run and support model, outstanding questions.\\
Architecture & Design, integrations with tickets, SharePoint, and Graph, data flows. & Target
architecture; decisions on interfaces and retrieval scope.\\
Security architecture & Permissions, identities, retrieval scope, logging evidence; ticket data. & Issue
the red opinion. Specify a dedicated identity, controlled actions, permission-filtered retrieval, and
logs; route redesign or the read-only pilot back for review.\\
Tech readiness & Revised design after re-review; pilot and test results; runbooks. & Check conditions
against actual access and action tests; produce readiness evidence and unresolved items.\\
General & Gate opinions, residual risks with business owners, budget, schedule, mandate. & Record the
decision for the defined scope and its conditions, and the sponsors' commitment. Generating the dossier
does not supply that authorization.\\
\bottomrule
\end{tabularx}
\end{table}

Automation can assemble the dossier, compare the declared read-only scope with actual permissions, and
link each condition to a test and a responsible owner. Difficult identity defects contribute to
$\kappa$; assistance with their resolution determines $\eta$; checking corrected conditions
contributes to $\mu$; repeated failed reviews contribute to rework; maintaining permission tests
contributes to $B_A$. Corrected evidence and authorized decisions, not generated documentation, clear
the red opinion.

The agent being evaluated and the system executing the DGF are separate principals, with separate
identities and authorities. A project cannot approve itself by generating its own evidence. A fully
automated arrangement preserves separation of duties through independent identities, evidence sources,
execution restrictions, and authorization policies; it need not identify separation of duties with two
human bodies. The sponsors' commitment at the general gate is deliberately not removed from the thesis:
it becomes the exercise of a standing mandate (Section~\ref{sec:mandate}). Reusable patterns and
testable environments make Build the most plausible first route to close.
